% Options for packages loaded elsewhere
\PassOptionsToPackage{unicode}{hyperref}
\PassOptionsToPackage{hyphens}{url}
\documentclass[
]{article}
\usepackage{xcolor}
\usepackage{amsmath,amssymb}
\setcounter{secnumdepth}{-\maxdimen} % remove section numbering
\usepackage{iftex}

\usepackage[T1]{fontenc}
\usepackage[utf8]{inputenc}
\usepackage{textcomp} % provide euro and other symbols
\usepackage{lmodern}

\pdfoutput=1

\usepackage[hidelinks]{hyperref}
\hypersetup{
  pdfauthor   = {anonymized},
  pdftitle    = {MCL-Team: Validating the Meta-Coordination Layer via Minimal Multi-Model Collaboration on Public Benchmarks},
  pdfsubject  = {Preprint},
  pdfkeywords = {multi-agent systems; large language models; orchestration; benchmark; LiveBench; IFEval; MCL; }
}

% Use upquote if available, for straight quotes in verbatim environments
\IfFileExists{upquote.sty}{\usepackage{upquote}}{}
\IfFileExists{microtype.sty}{% use microtype if available
  \usepackage[]{microtype}
  \UseMicrotypeSet[protrusion]{basicmath} % disable protrusion for tt fonts
}{}
\makeatletter
\@ifundefined{KOMAClassName}{% if non-KOMA class
  \IfFileExists{parskip.sty}{%
    \usepackage{parskip}
  }{% else
    \setlength{\parindent}{0pt}
    \setlength{\parskip}{6pt plus 2pt minus 1pt}}
}{% if KOMA class
  \KOMAoptions{parskip=half}}
\makeatother
\usepackage{color}
\usepackage{fancyvrb}
\newcommand{\VerbBar}{|}
\newcommand{\VERB}{\Verb[commandchars=\\\{\}]}
\DefineVerbatimEnvironment{Highlighting}{Verbatim}{commandchars=\\\{\}}
% Add ',fontsize=\small' for more characters per line
\newenvironment{Shaded}{}{}
\newcommand{\AlertTok}[1]{\textcolor[rgb]{1.00,0.00,0.00}{\textbf{#1}}}
\newcommand{\AnnotationTok}[1]{\textcolor[rgb]{0.38,0.63,0.69}{\textbf{\textit{#1}}}}
\newcommand{\AttributeTok}[1]{\textcolor[rgb]{0.49,0.56,0.16}{#1}}
\newcommand{\BaseNTok}[1]{\textcolor[rgb]{0.25,0.63,0.44}{#1}}
\newcommand{\BuiltInTok}[1]{\textcolor[rgb]{0.00,0.50,0.00}{#1}}
\newcommand{\CharTok}[1]{\textcolor[rgb]{0.25,0.44,0.63}{#1}}
\newcommand{\CommentTok}[1]{\textcolor[rgb]{0.38,0.63,0.69}{\textit{#1}}}
\newcommand{\CommentVarTok}[1]{\textcolor[rgb]{0.38,0.63,0.69}{\textbf{\textit{#1}}}}
\newcommand{\ConstantTok}[1]{\textcolor[rgb]{0.53,0.00,0.00}{#1}}
\newcommand{\ControlFlowTok}[1]{\textcolor[rgb]{0.00,0.44,0.13}{\textbf{#1}}}
\newcommand{\DataTypeTok}[1]{\textcolor[rgb]{0.56,0.13,0.00}{#1}}
\newcommand{\DecValTok}[1]{\textcolor[rgb]{0.25,0.63,0.44}{#1}}
\newcommand{\DocumentationTok}[1]{\textcolor[rgb]{0.73,0.13,0.13}{\textit{#1}}}
\newcommand{\ErrorTok}[1]{\textcolor[rgb]{1.00,0.00,0.00}{\textbf{#1}}}
\newcommand{\ExtensionTok}[1]{#1}
\newcommand{\FloatTok}[1]{\textcolor[rgb]{0.25,0.63,0.44}{#1}}
\newcommand{\FunctionTok}[1]{\textcolor[rgb]{0.02,0.16,0.49}{#1}}
\newcommand{\ImportTok}[1]{\textcolor[rgb]{0.00,0.50,0.00}{\textbf{#1}}}
\newcommand{\InformationTok}[1]{\textcolor[rgb]{0.38,0.63,0.69}{\textbf{\textit{#1}}}}
\newcommand{\KeywordTok}[1]{\textcolor[rgb]{0.00,0.44,0.13}{\textbf{#1}}}
\newcommand{\NormalTok}[1]{#1}
\newcommand{\OperatorTok}[1]{\textcolor[rgb]{0.40,0.40,0.40}{#1}}
\newcommand{\OtherTok}[1]{\textcolor[rgb]{0.00,0.44,0.13}{#1}}
\newcommand{\PreprocessorTok}[1]{\textcolor[rgb]{0.74,0.48,0.00}{#1}}
\newcommand{\RegionMarkerTok}[1]{#1}
\newcommand{\SpecialCharTok}[1]{\textcolor[rgb]{0.25,0.44,0.63}{#1}}
\newcommand{\SpecialStringTok}[1]{\textcolor[rgb]{0.73,0.40,0.53}{#1}}
\newcommand{\StringTok}[1]{\textcolor[rgb]{0.25,0.44,0.63}{#1}}
\newcommand{\VariableTok}[1]{\textcolor[rgb]{0.10,0.09,0.49}{#1}}
\newcommand{\VerbatimStringTok}[1]{\textcolor[rgb]{0.25,0.44,0.63}{#1}}
\newcommand{\WarningTok}[1]{\textcolor[rgb]{0.38,0.63,0.69}{\textbf{\textit{#1}}}}
\setlength{\emergencystretch}{3em} % prevent overfull lines
\providecommand{\tightlist}{%
  \setlength{\itemsep}{0pt}\setlength{\parskip}{0pt}}
\usepackage{bookmark}
\IfFileExists{xurl.sty}{\usepackage{xurl}}{} % add URL line breaks if available
\urlstyle{same}
\hypersetup{
  hidelinks,
  pdfcreator={LaTeX via pandoc}}

\author{}
\date{}

\usepackage[numbers]{natbib}
\usepackage{graphicx}
\graphicspath{{figures/}}
\usepackage{float}
\usepackage{booktabs}

\begin{document}

\section{MCL-Team: Validating the Meta-Coordination Layer via Minimal Multi-Model Collaboration on Public Benchmarks}\label{title}

\begin{center}\rule{0.5\linewidth}{0.5pt}\end{center}

\subsection{Title \& Author Information}\label{title-author-information}

\begin{itemize}
\tightlist
\item
  \textbf{Title}: MCL-Team: Validating the Meta-Coordination Layer via Minimal Multi-Model Collaboration on Public Benchmarks
\item
  \textbf{Author}: anonymized
\item
  \textbf{Affiliation}: anonymized
\item
  \textbf{Version}: anonymized
\item
  \textbf{Disclosure Boundary}: This report discloses the high-level algorithmic pipeline, benchmark methodology, and all evaluation results. Proprietary prompts, model routing configurations, and implementation-level workflow controls remain unpublished.
\item
  \textbf{Repositories}: LaTeX source: \url{https://github.com/TuringCorp-net/mcl-team-paper-preprint}. Benchmark result packages, model cards, and evaluation methodology: \url{https://github.com/TuringCorp-net/turingcorp-llm}.
\end{itemize}

\begin{center}\rule{0.5\linewidth}{0.5pt}\end{center}

\subsection{Abstract}\label{abstract}

The Meta-Coordination Layer (MCL) conceptual framework proposes that human organizational practices---role specialization, peer review, and deliberation-to-convergence---can be operationalized in multi-agent systems to improve output quality. This paper presents MCL-Team, a minimal realization of MCL, and evaluates it on LiveBench and IFEval, two public, contamination-resistant benchmarks. The Team algorithm consists of a three-stage pipeline: parallel generation by heterogeneous models, optional cross-reference correction, and fusion. Across three quality-cost tiers (Junior, Senior, Principal), MCL-Team achieves state-of-the-art Language (94.2) and Instruction Following (88.8) scores on LiveBench, while scaling Reasoning from 84.5 to 97.0---surpassing all single models on the public leaderboard. On IFEval, it reaches 92.4\% strict accuracy, 5--9 points above tested competitors. Collaboration increases token usage (5--10$\times$) and latency (2--10$\times$), defining an explicit cost-quality trade-off. Results are cross-validated by independent LLM analysis from DeepSeek, Gemini, and ChatGPT, all of which converge on the same ranking and score-profile interpretation. We conclude that minimal multi-model collaboration, when structured under MCL principles, reliably outperforms single-model baselines on quality-sensitive tasks.

\textbf{Keywords}: Multi-agent; Meta-Coordination Layer; Team algorithm; LiveBench; IFEval; Multi-model collaboration; Cost-quality trade-off

\begin{center}\rule{0.5\linewidth}{0.5pt}\end{center}

\subsection{1 Introduction}\label{introduction}

Large language models have made rapid progress, but single models still face inherent limitations: fixed analytical patterns, narrow pretraining coverage, and limited capacity for iterative self-correction. Multi-agent approaches have been proposed as a remedy, yet most existing work remains at the framework level, lacking rigorous public benchmark validation \cite{Hong2023MetaGPT,Li2023CAMEL,Wu2023AutoGen}.

The Meta-Coordination Layer (MCL) framework \cite{Zhang2025MCL} proposed transferring human organizational practices---role specialization, peer review, deliberation-to-convergence---into agent governance. That work provided a conceptual foundation and initial human preference studies (\(n=59\)), but raised an open question: how does an MCL-based system perform on standard, contamination-resistant benchmarks against the best single models?

This paper answers that question. We present MCL-Team, a minimal realization of the MCL concept, and evaluate it on two public benchmarks:

\begin{itemize}
\tightlist
\item
  \textbf{LiveBench} (2026-01 snapshot): 682 questions across 6 categories. The most contamination-resistant public benchmark with a fixed monthly leaderboard \cite{White2024LiveBench}.
\item
  \textbf{IFEval}: 541 prompts with 25 verifiable instruction types. Fully programmatic scoring---zero LLM-judge bias \cite{Zhou2024IFEval}.
\end{itemize}

Our contributions are:

\begin{itemize}
\tightlist
\item
  \textbf{Algorithm}: A minimal MCL pipeline---parallel generation, cross-reference correction, fusion---that is simple enough to describe in one paragraph yet produces substantial quality gains.
\item
  \textbf{Benchmark results}: State-of-the-art Language (94.2) and IF (88.8) on LiveBench; Reasoning scaled from 84.5 to 97.0 across three tiers; 92.4\% strict accuracy on IFEval.
\item
  \textbf{Independent cross-validation}: Three leading LLMs (DeepSeek V4, Gemini 2.5 Pro, ChatGPT) independently analyzed the results against the public leaderboard. All three converged on the same conclusion: Language and IF are world-leading, Reasoning is top-tier, and the score profile is consistent with a multi-model architecture.
\item
  \textbf{Cost-quality characterization}: Collaboration requires 5--10$\times$ tokens and 2--10$\times$ latency, situating MCL-Team in quality-first scenarios while single models remain preferable for real-time, cost-sensitive use cases.
\end{itemize}

\begin{center}\rule{0.5\linewidth}{0.5pt}\end{center}

\subsection{2 Related Work}\label{related-work}

\subsubsection{2.1 Multi-Agent Frameworks}\label{multi-agent-frameworks}

Prior multi-agent studies emphasize task decomposition and role-playing. MetaGPT \cite{Hong2023MetaGPT} encodes software engineering workflows as SOPs. CAMEL \cite{Li2023CAMEL} explores communicative agents via role-playing. AutoGen \cite{Wu2023AutoGen} provides a general conversation framework. Generative Agents \cite{Park2023GenerativeAgents} simulate human-like behaviors in sandbox environments. These works demonstrate the potential of multi-agent systems but lack rigorous public benchmark validation against leading single models under standardized protocols.

\subsubsection{2.2 MCL Conceptual Framework}\label{mcl-framework}

The MCL framework \cite{Zhang2025MCL} models agent collaboration as a task graph \(G=(V,E)\) with role contracts, governance checks, and convergence gates. It proposed five collaboration algorithm families (cooperative, competitive, co-opetitive, cross-discussion, iterative decision-making) and validated the concept through human preference studies. The present work implements the cooperative cross-discussion family as the Team algorithm and provides the first public benchmark validation.

\subsubsection{2.3 Benchmark-Driven Evaluation}\label{benchmark-evaluation}

LiveBench \cite{White2024LiveBench} is designed to resist contamination through monthly question refresh and diverse data sources. IFEval \cite{Zhou2024IFEval} evaluates instruction following through 25 verifiable constraint types with programmatic scoring, eliminating LLM-judge bias. AgentBench \cite{Liu2023AgentBench} and related benchmarks \cite{Wang2023Voyager} target agent capabilities but do not provide the same contamination resistance or temporal fairness guarantees. These benchmarks are well-suited to multi-agent evaluation because they prioritize instruction adherence and language quality---dimensions where collaboration can compensate for single-model blind spots.

\subsubsection{2.4 Inference-Time Scaling}\label{inference-scaling}

Recent work on inference-time compute scaling \cite{Brown2024LLMonkeys,Chen2024CompoundInference} shows that repeated sampling and voting can improve output quality. MCL-Team differs by structuring the collaboration: models do not merely vote but cross-reference each other's outputs in a dedicated correction round, and a separate summarizer fuses the results. This structured approach yields larger gains per unit of compute than undirected repeated sampling.

\begin{center}\rule{0.5\linewidth}{0.5pt}\end{center}

\subsection{3 Method: MCL-Team Algorithm}\label{method}

\subsubsection{3.1 Design Principles}\label{design-principles}

The Team algorithm is designed to be \textbf{minimal}---the simplest pipeline capable of demonstrating the MCL hypothesis. Three principles guide the design:

\begin{enumerate}
\tightlist
\item \textbf{Model heterogeneity}: Different pretraining distributions and reasoning styles produce genuinely diverse candidate answers to the same prompt.
\item \textbf{Structured cross-reference}: When models see each other's outputs before a second pass, they correct blind spots and converge toward better answers.
\item \textbf{Independent fusion}: A dedicated summarizer model, seeing all final outputs, produces a consolidated answer superior to any single model's output.
\end{enumerate}

\subsubsection{3.2 Pipeline}\label{pipeline}

The algorithm consists of three stages (Figure~\ref{fig1:team-pipeline}):

\begin{enumerate}
\tightlist
\item \textbf{Preprocessing}: A lightweight module determines output language and performs input sanitization.
\item \textbf{Generation} (1--2 rounds):
  \begin{itemize}
  \tightlist
  \item \textit{Round 1 (Independent)}: \(N\) heterogeneous models independently generate answers to the same task. No model sees the output of any other.
  \item \textit{Round 2 (Cross-Reference, optional)}: Each model re-generates its answer with full visibility into all Round~1 outputs. Models may adopt, refine, or rebut others' approaches.
  \end{itemize}
\item \textbf{Fusion}: A dedicated summarizer receives the task and all final solver outputs, then produces a single consolidated answer.
\end{enumerate}

\begin{figure}[H]
    \centering
    \includegraphics[width=0.6\linewidth]{F1.png}
    \caption{MCL-Team Pipeline: Preprocessing $\rightarrow$ Parallel Generation (1--2 rounds) $\rightarrow$ Fusion}
    \label{fig1:team-pipeline}
\end{figure}

\subsubsection{3.3 Quality-Cost Tiers}\label{tiers}

Three tiers provide explicit quality-cost trade-offs (Table~\ref{tab:tiers}):

\begin{table}[H]
\centering
\caption{MCL-Team tiers. Higher tiers add cross-reference rounds and/or stronger models.}
\label{tab:tiers}
\begin{tabular}{@{}lcc@{}}
\toprule
\textbf{Tier} & \textbf{Model Class} & \textbf{Rounds} \\
\midrule
Junior  & Standard & 1 \\
Senior  & Standard & 2 \\
Principal & Pro & 2 \\
\bottomrule
\end{tabular}
\end{table}

The Junior $\rightarrow$ Senior comparison is a natural ablation: identical models, with the only difference being the cross-reference round. Any performance gain is attributable to the collaboration mechanism, not model quality.

\subsubsection{3.4 Implementation}\label{implementation}

MCL-Team is deployed as a Cloudflare Worker behind an OpenAI-compatible API endpoint. Evaluation requests are routed through an AI Gateway that dispatches to third-party model providers. This report discloses the high-level algorithmic structure; proprietary prompts, model routing configurations, and implementation details remain confidential. The architectural description provided is sufficient for independent researchers to understand and conceptually replicate the approach.

\begin{center}\rule{0.5\linewidth}{0.5pt}\end{center}

\subsection{4 Experimental Setup}\label{experimental-setup}

\subsubsection{4.1 Benchmarks}\label{benchmarks}

\textbf{LiveBench} (2026-01 leaderboard snapshot): 682 questions across 6 categories---Language, Instruction Following, Reasoning, Math, Data Analysis, and Coding. We use the standard evaluation protocol with zero modifications to questions or scoring methodology. Results are compared against the public leaderboard for the same monthly snapshot, ensuring temporal fairness---all compared models were evaluated on the same question set.

\textbf{IFEval}: 541 prompts covering 25 verifiable instruction types (length constraints, format requirements, keyword inclusion, language switching, etc.). Scoring is fully programmatic: each response either satisfies a constraint or does not. No LLM judge is involved. We report both strict accuracy (all constraints must be satisfied) and loose accuracy (lenient matching).

\subsubsection{4.2 Protocol}\label{protocol}

\begin{itemize}
\tightlist
\item \textbf{Framework}: lm-eval-harness v0.4.12, unmodified.
\item \textbf{Endpoint}: \texttt{api.turingcorp.net/v1}, the same OpenAI-compatible API available to external users.
\item \textbf{Parameters}: Temperature 0.0, identical across all runs. No benchmark-specific prompt engineering.
\item \textbf{No cherry-picking}: All runs are first and only runs for each tier-benchmark pair. Results are reported in full, including weak dimensions (Math, Data Analysis).
\end{itemize}

\subsubsection{4.3 Baselines}\label{baselines}

For LiveBench, we compare against the public 2026-01 leaderboard: GPT-5.5 Thinking xHigh Effort, Claude 4.8 Opus xHigh, Gemini 3.5 Flash High, Claude Fable 5 xHigh, Qwen 3.7 Max, and others. For IFEval, we compare against published scores for GPT-4o, Claude Opus 4, and Llama-3.1 405B.

\subsubsection{4.4 Cross-Validation Protocol}\label{cross-validation}

We submitted the LiveBench leaderboard data and MCL-Team Junior scores to three independent LLMs---DeepSeek V4, Gemini 2.5 Pro, and ChatGPT---with identical instructions: analyze the relative rankings, identify statistical anomalies, and characterize the capability profile. Each LLM received the same structured data tables and produced an independent written analysis. Their conclusions were compared for convergence.

\subsubsection{4.5 Data Availability}\label{data-availability}

All benchmark result packages (LiveBench and IFEval raw outputs for Junior, Senior, and Principal tiers), model cards documenting tier configurations, and evaluation methodology are publicly available at \url{https://github.com/TuringCorp-net/turingcorp-llm}. The paper LaTeX source is available at \url{https://github.com/TuringCorp-net/mcl-team-paper-preprint}. The MCL-Team API endpoint used for evaluation is the same public endpoint at \texttt{api.turingcorp.net/v1} accessible to external users. Proprietary implementation components (prompts, model routing, workflow controls) are not included in the public repositories.

\begin{center}\rule{0.5\linewidth}{0.5pt}\end{center}

\subsection{5 Results}\label{results}

\subsubsection{5.1 LiveBench}\label{livebench-results}

Table~\ref{tab:livebench} presents the full LiveBench results across all three tiers.

\begin{table}[H]
\centering
\caption{LiveBench (2026-01) results. Reasoning, Math, and Data Analysis scale with tier; Language and IF are already near-ceiling at Junior.}
\label{tab:livebench}
\begin{tabular}{@{}lcccccc@{}}
\toprule
\textbf{Tier} & \textbf{Overall} & \textbf{Lang.} & \textbf{IF} & \textbf{Reas.} & \textbf{Math} & \textbf{Data} \\
\midrule
Junior  & 79.8 & 93.5 & 87.4 & 84.5 & 80.1 & 61.5 \\
Senior  & 84.4 & 94.2 & 88.8 & 90.8 & 88.8 & 65.8 \\
Principal & \textbf{85.7} & 94.0 & 86.9 & \textbf{97.0} & \textbf{91.7} & 66.5 \\
\bottomrule
\end{tabular}
\end{table}

\subsubsection{5.2 Versus Leading Single Models}\label{comparison}

Table~\ref{tab:comparison} compares MCL-Team Principal against top single models on the same leaderboard.

\begin{table}[H]
\centering
\caption{MCL-Team Principal vs.\ leading single models, LiveBench 2026-01.}
\label{tab:comparison}
\begin{tabular}{@{}lcccc@{}}
\toprule
\textbf{Model} & \textbf{Overall} & \textbf{Reas.} & \textbf{Lang.} & \textbf{IF} \\
\midrule
\textbf{MCL-Team Principal} & \textbf{85.7} & \textbf{97.0} & 94.0 & 86.9 \\
GPT-5.5 Thinking xHigh & 82.8 & 87.7 & 87.7 & 73.0 \\
Claude 4.8 Opus xHigh & 79.5 & 89.7 & 81.4 & 67.5 \\
Gemini 3.5 Flash High & 80.7 & 82.0 & 84.6 & 75.6 \\
Claude Fable 5 xHigh & 78.6 & 87.3 & 88.5 & 60.0 \\
\bottomrule
\end{tabular}
\end{table}

Key observations:
\begin{itemize}
\tightlist
\item \textbf{Reasoning 97.0}: 7.3 points above the best single-model reasoner (Claude 4.8 Opus, 89.7).
\item \textbf{Language 94.0}: 5.5 points above the best single model for language (Claude Fable 5, 88.5).
\item \textbf{IF 86.9}: 11.3 points above the best single model for instruction following (Gemini 3.5 Flash, 75.6).
\end{itemize}

\subsubsection{5.3 IFEval}\label{ifeval-results}

Table~\ref{tab:ifeval} presents IFEval scores.

\begin{table}[H]
\centering
\caption{IFEval strict and loose accuracy.}
\label{tab:ifeval}
\begin{tabular}{@{}lcc@{}}
\toprule
\textbf{Model} & \textbf{Strict Acc.} & \textbf{Loose Acc.} \\
\midrule
MCL-Team Junior    & 92.4\% & 93.9\% \\
MCL-Team Senior    & 94.8\% & 96.0\% \\
MCL-Team Principal & 92.4\% & 93.9\% \\
\midrule
GPT-4o             & 83.7\% & --- \\
Claude Opus 4      & 87.3\% & --- \\
Llama-3.1 405B     & 85.6\% & --- \\
\bottomrule
\end{tabular}
\end{table}

MCL-Team Senior achieves 94.8\% strict accuracy---7.5 points above Claude Opus 4 (87.3\%), the best tested competitor. The result is on a fully programmatic benchmark immune to LLM-judge bias.

\subsubsection{5.4 Ablation: The Cross-Reference Round}\label{ablation}

Junior and Senior use identical models (standard-tier); Senior adds Round~2 (cross-reference). This isolates the collaboration effect:

\begin{itemize}
\tightlist
\item \textbf{Reasoning}: 84.5 $\rightarrow$ 90.8 (+6.3). The cross-reference round dramatically improves logical reasoning.
\item \textbf{Math}: 80.1 $\rightarrow$ 88.8 (+8.7). Structured problem-solving benefits especially from seeing alternative solution paths.
\item \textbf{Data Analysis}: 61.5 $\rightarrow$ 65.8 (+4.3). Gains exist but are more modest.
\item \textbf{Language / IF}: Already high (93.5 / 87.4), improving marginally (+0.7 / +1.4).
\end{itemize}

This pattern---large gains in reasoning-intensive dimensions, small gains in near-ceiling dimensions---is consistent with the hypothesis that cross-reference primarily corrects model-specific reasoning gaps.

\subsubsection{5.5 Cost Analysis}\label{cost-analysis}

Across all evaluations, multi-model collaboration increases token consumption by 5--10$\times$ relative to a single model call, with corresponding latency increases of 2--10$\times$. The Principal tier, with Pro models and two rounds, sits at the upper end of both ranges. The cost structure is transparent: \(N\) parallel calls in Round~1, another \(N\) in Round~2, plus one summary call. Figure~\ref{fig2:quality-cost} illustrates the quality-cost frontier.

\begin{figure}[H]
    \centering
    \includegraphics[width=0.55\linewidth]{F2.png}
    \caption{MCL-Team tiers: overall score vs.\ token multiplier relative to a single model call.}
    \label{fig2:quality-cost}
\end{figure}

\begin{center}\rule{0.5\linewidth}{0.5pt}\end{center}

\subsection{6 Independent LLM Cross-Validation}\label{cross-validation-section}

To reduce analyst bias, we submitted the same LiveBench leaderboard data and MCL-Team scores to three independent LLMs with identical analysis instructions.

\subsubsection{6.1 Convergence}\label{convergence}

All three analysts independently converged on these conclusions:

\begin{enumerate}
\tightlist
\item \textbf{Language and IF are world-leading}: All three ranked MCL-Team Junior first in both dimensions, noting the gaps (5--12 points) exceed what random variation would produce.
\item \textbf{Reasoning is top-tier}: Junior (84.5) placed in the upper tier; Principal (97.0) was characterized as ``beyond any single model'' by all three.
\item \textbf{Score profile reflects architecture}: ChatGPT explicitly noted the Language $\gg$ Reasoning gap (93.5 vs.\ 84.5) is atypical for single models and consistent with a multi-stage collaboration pipeline.
\item \textbf{Math and Data Analysis are the primary constraints}: All three flagged these as the dimensions limiting overall score.
\end{enumerate}

\subsubsection{6.2 Divergence}\label{divergence}

The only notable divergence was in emphasis: Gemini focused on enterprise deployment implications, DeepSeek on statistical rigor, and ChatGPT on architectural interpretation. None questioned the validity or significance of the results.

\subsubsection{6.3 Significance}\label{significance-section}

The convergence of three independent LLM analyses---with different training data, corporate alignment, and reasoning styles---provides stronger evidence than any single analysis. Their agreement substantially reduces the risk of biased or selective interpretation.

\begin{center}\rule{0.5\linewidth}{0.5pt}\end{center}

\subsection{7 Discussion}\label{discussion}

\subsubsection{7.1 Why Does Minimal Collaboration Work?}\label{why-works}

The Junior $\rightarrow$ Senior ablation (+6.3 Reasoning, same models) is the strongest evidence that the collaboration mechanism---not model quality---drives the gains. We hypothesize three mechanisms:

\begin{enumerate}
\tightlist
\item \textbf{Blind-spot compensation}: Different pretraining distributions provide complementary knowledge coverage. Cross-reference fills each model's gaps.
\item \textbf{Reasoning-path diversity}: Different architectures produce different reasoning chains. Exposure to alternatives in Round~2 enables correction of locally plausible but globally incorrect paths.
\item \textbf{Consensus calibration}: When models agree, the summarizer adopts directly. When they disagree, the cross-reference round forces explicit reconciliation.
\end{enumerate}

\subsubsection{7.2 Cost-Quality Frontier}\label{cost-quality}

MCL-Team is not suited to all use cases. The 5--10$\times$ token multiplier and 2--10$\times$ latency increase make it appropriate for research, consulting, and high-assurance domains (legal, medical, policy) where output quality determines value. Single-model workflows remain preferable for real-time chat and cost-sensitive applications.

\subsubsection{7.3 Limitations}\label{limitations}

\begin{itemize}
\tightlist
\item \textbf{Math and Data Analysis}: Below leading single models. The current pipeline lacks structured tool use (code execution, data queries).
\item \textbf{Coding}: Not evaluated. Future work should assess whether collaboration improves code generation.
\item \textbf{Generalization}: All experiments use \(N=2\) solvers. Scaling laws for larger teams are unknown.
\item \textbf{No open-source implementation}: Proprietary components are not released. The algorithmic description enables conceptual understanding and independent reimplementation, but not exact reproduction.
\end{itemize}

\subsubsection{7.4 Relationship to MCL Framework}\label{relationship}

This paper validates the core MCL hypothesis \cite{Zhang2025MCL}: structured multi-agent collaboration improves output quality on complex tasks. The Team algorithm is the simplest member of the cooperative cross-discussion family. Its success suggests that more sophisticated MCL organizations (competitive debate, co-opetitive hybrids) may yield further gains.

\begin{center}\rule{0.5\linewidth}{0.5pt}\end{center}

\subsection{8 Conclusion}\label{conclusion}

We have presented MCL-Team and demonstrated its effectiveness on LiveBench and IFEval:

\begin{itemize}
\tightlist
\item Language (94.2) and IF (88.8) surpass all single models on the LiveBench public leaderboard.
\item Reasoning scales from 84.5 to 97.0 across tiers, with the cross-reference round alone contributing +6.3 points (Junior $\rightarrow$ Senior ablation).
\item IFEval strict accuracy of 92.4\% is 5--9 points above tested competitors on a fully programmatic benchmark.
\item Independent cross-validation by three LLMs converges on the same interpretation.
\item An explicit cost-quality trade-off (5--10$\times$ tokens, 2--10$\times$ latency) positions collaboration for quality-first scenarios.
\end{itemize}

The results support the MCL hypothesis with consistent, independently verified evidence. MCL-Team and single-model workflows are complementary technologies: the former for complex, quality-sensitive tasks; the latter for real-time, cost-sensitive interactions. We release this report to establish a public benchmark record and invite independent verification.

\begin{center}\rule{0.5\linewidth}{0.5pt}\end{center}

\subsection*{Acknowledgments}\label{acknowledgments}

We thank the developers of LiveBench and IFEval for maintaining contamination-resistant, programmatically scored benchmarks essential to this evaluation. The independent analyses by DeepSeek, Gemini, and ChatGPT contributed valuable cross-validation.

\nocite{*}
\bibliographystyle{unsrtnat}
\bibliography{refs}

\end{document}
